\begin{center}
\textsc{Gauss} an \textsc{Bessel}.\quad Göttingen, 9.~April 1830.
\end{center}

\par\noindent\hspace*{\parindent}.\ .\ .\ .\ .\ .\ Noch viel mehr Zeit haben mir seit Mai 1829 die trigonometrischen Messungen geraubt, wenn ich gleich keinen unmittelbaren Antheil an den Geschäften im Felde diesmal genommen habe.
Noch diese Stunde bin ich nicht ganz (obwohl Gott Lob beinahe) mit Verarbeitung der vorigjährigen Messungen fertig, wobei ich jeder Hülfe entbehre.
Doch habe ich dabei viel Freude über die zu meiner grössten Zufriedenheit ausgefallene Art gehabt, wie mein Sohn seinen Antheil an diesen Geschäften ausgeführt hat.
Er hat ganz allein ein grosses Dreiecksnetz von der Weser bis zur holländischen Grenze (über das Osnabrücksche, auch mit mehrmaliger Berührung preussischen Gebiets, wobei ihm von Seiten der dortigen Behörden sehr liberal Vorschub geleistet ist) geführt, wodurch ausser den Hauptpunkten noch gegen 250 andere festgelegt sind.
Die --- von einem andern Officier --- im Fürstenthum Hildesheim gemachten trigonometrischen Messungen sind gleichfalls so gut wie vollendet, so dass alle Messtischblätter (53, wovon etwa ein Viertel bereits aufgenommen ist) mit festen Punkten versehen werden können.\ .\ .\ .\ .\ .\ .

\begin{center}\rule{.24\linewidth}{.4pt}\end{center}
% The printer's signature is visible in the Google witness, physical leaf 372.
\par\smallskip\noindent\hfill{\footnotesize 46\textsuperscript{*}}\par
